\section{The arithmetic target and its proof discipline}
\label{sec:foundations}

\subsection{The theory used throughout}

The object theory is first-order intuitionistic arithmetic with natural
numbers as its only sort. Its language has \(0,S\) and a symbol for every
fixed primitive recursive construction. The symbols are generated from
zero, successor and projections by composition and primitive recursion.
For a recursion symbol the defining equations are
\[
 r(0,\boldsymbol x)=g(\boldsymbol x),\qquad
 r(n+1,\boldsymbol x)=h(n,r(n,\boldsymbol x),\boldsymbol x).
\]
We include the successor axioms \(Sx\ne0\) and
\(Sx=Sy\to x=y\), equality reflexivity and substitution, and ordinary
intuitionistic first-order natural deduction. In particular implication,
conjunction, disjunction, falsity and both quantifiers have their usual
introduction and elimination rules, with the standard freshness conditions
on quantifier introduction and existential elimination. No excluded-middle
schema, choice principle or Markov principle is added.

Induction is restricted to quantifier-free formulas, with parameters:
\[
 A(0,\boldsymbol x),\quad
 \forall n\bigl(A(n,\boldsymbol x)\to A(n+1,\boldsymbol x)\bigr)
 \quad\Longrightarrow\quad \forall n\,A(n,\boldsymbol x).
\]
This is the ordinary \ipra\ presentation formalized by
\texttt{PRAAx}, \texttt{PTm}, and \texttt{PR} in ConciseIPRA
\cite{repoCore}. The native concise calculus is connected to it by an
existing syntactic bi-interpretation. We work first in this explicit
arithmetic presentation, and transport a derivation only after it exists.

The induction restriction does not prohibit formulas with quantifiers from
appearing as assumptions or conclusions. For example, a hypothesis
\(\forall n\,C(n)\) may be instantiated at each of finitely many indices in
a calculation. What must be audited is the formula on which induction is
performed.

\subsection{Finite objects, total operations, and decidability}

An integer is a signed natural-number code; a rational is a pair
\((z,d)\) with \(d>0\). We do not need quotient objects: equality is
\((z,d)\simeq(w,e)\) precisely when \(ze=wd\). Addition, multiplication,
comparison and inversion with a nonzero-numerator guard are explicit
integer calculations. A gcd normalization is optional. All operations
receive harmless default values on malformed input; the theorem carries
the validity guards.

Finite tuples, coefficient arrays, matrices and execution histories are
also natural-number codes. One possible list code is empty list \(0\)
and
\[
 \operatorname{cons}(a,t)=1+2^a(2t+1).
\]
Recovering its head and tail uses bounded tests on a code; the tail is
smaller than a nonempty valid code. Repeated decoding can therefore be
implemented with a loop bounded by that code. This very inefficient coding
is nevertheless primitive recursive. A practical implementation may use
a more efficient pairing scheme; no argument depends on efficiency.

A polynomial is a finite coefficient list or exponent-coefficient table.
Multiplication is a bounded convolution. A determinant is the signed sum
over finitely many permutations, with at most \(n!\) summands. A rational
matrix can be row-reduced by scanning bounded lists of columns and rows;
each pivot is selected by a decidable nonzero test. These descriptions
establish that the operations are primitive recursive. Their algebraic
correctness still needs proofs, given by the finite invariants described
below. A computation model is not a correctness argument.

\begin{lemma}[Booleanization, as a schema]
\label{lem:booleanization}
For every fixed bounded arithmetic formula \(B(\boldsymbol x)\), built
from fixed PR terms and bounded number quantifiers, there is a fixed PR
function \(b_B\) with values in \(\{0,1\}\), and an \ipra\ derivation
\[
 B(\boldsymbol x)\ \longleftrightarrow\ b_B(\boldsymbol x)=1.
\]
Consequently induction on bounded formulas is a derived use of
quantifier-free induction.
\end{lemma}
\begin{proof}
Equality and order have their ordinary PR Boolean tests. The successor
axioms and quantifier-free induction prove their correctness. Boolean
complement, product and maximum implement negation, conjunction and
disjunction. The Boolean implication test returns one exactly when the
antecedent test is zero or the consequent test is one. Its equivalence to
intuitionistic implication uses the already proved decidability of the
antecedent; it does not decide arbitrary propositions.

For a bounded universal quantifier use a prefix product
\(p(0)=1\), \(p(t+1)=p(t)c(t)\), where \(c(t)\in\{0,1\}\).
Quantifier-free induction proves \(p(t)\in\{0,1\}\) and
\(p(t+1)\le p(t)\). Fix an arbitrary index \(i\). Induction on the
quantifier-free implication
\(i<t\land p(t)=1\to c(i)=1\) proves the forward direction of
\(p(b)=1\leftrightarrow\forall i<b\,c(i)=1\).
For the reverse direction, keep the bounded universal assertion as a
fixed assumption and induct on \(t\le b\to p(t)=1\). In the step
\(t<b\), instantiate that assumption at \(t\).

For bounded existence scan all positions below the bound, retaining the
first success and using the current scan length as a failure sentinel
(so the final sentinel is the final bound). Induct
on the conjunction of: the stored value is at most the current length;
if it is smaller it is a successful witness; and the prefix product of
the failure tests before it is one. These are quantifier-free tests.
A supplied successful position precludes the sentinel, and the stored
position then witnesses the existential assertion. This proves both
directions without double-negation elimination.

Structural induction in the metatheory on the fixed formula supplies
the test and its derivation. Replacing a bounded induction formula by
the equation \(b_B(n,\boldsymbol x)=1\) completes the last assertion.
\end{proof}

This proof also gives bounded minimization and finite choice. In particular,
from a proof of
\(\forall i<n\,\exists y<B(x,i)\,R(x,i,y)\), with fixed PR bounds and
tests, a PR scan produces the finite list of least witnesses. There is no
corresponding inference from unbounded existence without an additional
bound or a separately proved PR constructor.

\subsection{A precise finite-loop rule}

Let the fixed PR run be
\[
 S(x,0)=\operatorname{init}(x),\qquad
 S(x,t+1)=\operatorname{step}(x,t,S(x,t)).
\]
Let \(P(x)\), \(I(x,N,t,s)\), and \(Q(x,N,s)\) be bounded predicates.
Suppose there are already \ipra\ proofs of
\begin{align*}
 P(x)&\to I(x,N,0,\operatorname{init}(x)),\\
 t<N\land I(x,N,t,s)&\to
 I(x,N,t+1,\operatorname{step}(x,t,s)),\\
 I(x,N,N,s)&\to Q(x,N,s).
\end{align*}
Then \ipra\ proves \(P(x)\to Q(x,N,S(x,N))\). Indeed, under the
fixed assumption \(P(x)\), induct on the Booleanization of
\[
              t\le N\ \longrightarrow\ I(x,N,t,S(x,t)).
\]
For the successor case with \(t+1\le N\), apply the step proof; outside
the guard the implication is immediate. Instantiating at \(N\) gives
the conclusion. This is the precise meaning of a finite induction audit
in later chapters.

The rule also accommodates a finite log of hypotheses queried from a
represented input. Include a Boolean conjunction saying that every
previously queried fact holds. A global regularity assumption can supply
the next fact by universal elimination, or the entire finite log can be
part of the theorem's input. The induction invariant remains a bounded
predicate about the coded state.

\begin{theorem}[Bounded descent with a reconstruction invariant]
\label{thm:rank}
Fix PR functions \(T,s_0,r\), a Boolean terminal test \(D\), and a
bounded validity predicate \(V\). Suppose \ipra\ proves that \(T\)
preserves \(V\), fixes terminal valid states, and satisfies
\[
 V(s)\land D(s)=0\quad\Longrightarrow\quad r(T(s))<r(s).
\]
If \(V(s_0)\), then the iterate \(T^{r(s_0)}(s_0)\) is terminal.
Every bounded invariant proved to be preserved by \(T\) holds there.
\end{theorem}
\begin{proof}
The bounded-loop rule gives validity of the iterates. Induct on
\[
 D(T^t(s_0))=1\quad\lor\quad r(T^t(s_0))+t\le r(s_0).
\]
If the terminal alternative holds, absorption preserves it. Otherwise
strict decrease of the natural-number rank preserves the second
alternative at \(t+1\). At \(t=r(s_0)\), a nonterminal state would
have rank zero. Its stipulated next step would have negative rank, a
contradiction. Hence the terminal test is one. A simultaneous bounded
invariant is proved by the same loop rule.
\end{proof}

This is background infrastructure, not a new Mordell--Weil argument.
The inspected Mordell--Weil packet already proves a more efficient
contracting-height descent once actual quotient representatives are
supplied. The theorem is useful here because a progress measure, a
coefficient precision and a finite row-reduction phase are all explicit
natural-number budgets. It also has a compiled Lean proof.

\subsection{What is uniform, and what is only a schema}

The degrees, matrix dimensions, integer moduli, coefficient lists and
precision bounds in the applications may be variable natural-number
inputs. Their algorithms use a fixed finite vocabulary of integer and
rational operations, together with explicitly bounded loops. They are
therefore genuinely uniform PR constructions.

An abstract coefficient ring is different. A theorem over every
\emph{PR-presented ring} is a schema: fix PR operations and supply or prove
their ring laws; then compile that instance. This does not introduce a
variable ranging over all total PR programs with a universal evaluator.
If such a PR evaluator \(U(e,x)\) existed, listing all unary PR functions,
then \(d(n)=U(n,n)+1\) would itself be PR and have an index \(e\).
Substitution of \(e\) gives \(U(e,e)=U(e,e)+1\), impossible. The
diagonal argument also explains why a generic evaluator hidden in an
apparently convenient Lean record would be a serious error.

Finally, a PR proof checker for arbitrary arithmetic proofs is not an
internal reflection principle. We never infer a formula from the single
assumption that an arbitrary coded proof passes a checker. We construct
derivations of the specific fixed formulas at issue. The object
composition lemmas in the archive transform actual \texttt{Prf PRAAx}
objects, with every input proof retained as a hypothesis.
